% ============================================================
% PAGE 6 — Module 2: Feature Fusion and Attention Mechanisms (Neck)
% ============================================================

\subsection{Module 2: Feature Fusion and Attention Mechanisms
(Neck)}

A backbone, however powerful, produces feature maps whose
representational strengths are inherently scale-dependent:
shallow stages preserve fine spatial detail but lack semantic
abstraction, while deep stages capture rich category-level
semantics at the expense of spatial resolution. The
\textit{neck} of a detection architecture exists precisely to
reconcile this dichotomy, constructing a unified multi-scale
feature pyramid in which every level carries both high
resolution \textit{and} strong semantics. In FusionDet the neck
comprises two tightly integrated components: a \textit{Bidirectional
Path Aggregation Network} (BiPANet) for structural feature
fusion and a \textit{Lightweight Channel Attention} (LCA)
module that recalibrates channel responses at every pyramid
level.

\subsubsection{Bidirectional Path Aggregation}

Classical Feature Pyramid Networks \cite{lin2017feature}
propagate information in a single top-down direction: the
deepest, most semantic feature map is progressively upsampled
and merged with shallower maps through lateral connections.
While effective, this unidirectional flow means that the
shallowest pyramid level receives strong semantic enrichment
but the deepest level gains little additional spatial detail.
The Path Aggregation Network (PANet) \cite{liu2018path}
addressed this asymmetry by appending a bottom-up pathway after
the top-down pass, allowing fine-grained spatial cues to flow
back toward the deeper levels.

FusionDet adopts and extends this bidirectional strategy. Let
$\{\mathbf{P}_{2},\mathbf{P}_{3},\mathbf{P}_{4},
\mathbf{P}_{5}\}$ denote the fused backbone outputs at strides
$\{4,8,16,32\}$ produced by the AFF module
(Section~\ref{sec:aff}). The top-down pass generates intermediate
representations $\{\mathbf{T}_{l}\}$ by recursive upsampling
and element-wise addition:
\begin{equation}
    \mathbf{T}_{l} =
    \phi^{1\times1}_{l}\!\bigl(\mathbf{P}_{l}\bigr)
    \;+\;
    \text{Up}_{2\times}\!\bigl(\mathbf{T}_{l+1}\bigr),
    \quad l = 4,3,2
    \label{eq:topdown}
\end{equation}
where $\phi^{1\times1}_{l}$ is a $1\times1$ convolution that
aligns channel dimensionality and $\text{Up}_{2\times}$ denotes
nearest-neighbour upsampling by a factor of two. The seed is
$\mathbf{T}_{5} = \phi^{1\times1}_{5}(\mathbf{P}_{5})$.

The subsequent bottom-up pass produces the final pyramid
features $\{\mathbf{N}_{l}\}$:
\begin{equation}
    \mathbf{N}_{l} =
    \text{DWSep}^{3\times3}_{l}\!\Bigl(
        \mathbf{T}_{l}
        \;+\;
        \text{Down}_{2\times}\!\bigl(\mathbf{N}_{l-1}\bigr)
    \Bigr),
    \quad l = 3,4,5
    \label{eq:bottomup}
\end{equation}
with $\mathbf{N}_{2} =
\text{DWSep}^{3\times3}_{2}(\mathbf{T}_{2})$ as the seed and
$\text{Down}_{2\times}$ implemented as a stride-2 depthwise
convolution rather than a pooling operation, preserving
learnable downsampling dynamics. Every $3\times3$ convolution in
both passes is realized as a depthwise-separable operation
\cite{howard2017mobilenets}, keeping the neck's total cost
below $0.8$~GFLOPs---approximately one-eighth of a vanilla
PANet with standard convolutions at the same channel width.

The bidirectional design ensures that $\mathbf{N}_{2}$ retains
the semantic boost injected during the top-down pass, while
$\mathbf{N}_{5}$ recovers spatial precision propagated upward
during the bottom-up pass. This symmetry is particularly
beneficial for small-object detection, where the shallowest
pyramid level must carry sufficient category-level information
to disambiguate visually similar but semantically distinct
targets \cite{liu2018path}.

\subsubsection{Lightweight Channel Attention (LCA)}

Even after bidirectional aggregation, not all channels within a
given pyramid level contribute equally to detection performance.
Some channels may encode background texture that is irrelevant
to the objects of interest, while others may carry precisely
the edge or color gradient information needed to resolve a
difficult instance. To allow the network to amplify informative
channels and suppress redundant ones, we apply a Lightweight
Channel Attention module after each pyramid output
$\mathbf{N}_{l}$.

Given $\mathbf{N}_{l} \in \mathbb{R}^{C \times H_l \times
W_l}$, LCA first aggregates spatial information into a channel
descriptor via global average pooling:
\begin{equation}
    z_{c}^{l} = \frac{1}{H_l \, W_l}
    \sum_{i=1}^{H_l}\sum_{j=1}^{W_l}
    \mathbf{N}_{l}(c,\,i,\,j)
    \label{eq:lca_gap}
\end{equation}
The descriptor vector
$\mathbf{z}^{l} \in \mathbb{R}^{C}$ is then passed through a
two-layer bottleneck with a reduction ratio $r$:
\begin{equation}
    \mathbf{a}^{l} = \sigma\!\Bigl(
        \mathbf{W}_{2}^{l}\;\text{SiLU}\!\bigl(
            \mathbf{W}_{1}^{l}\,\mathbf{z}^{l}
        \bigr)
    \Bigr)
    \label{eq:lca_gate}
\end{equation}
where $\mathbf{W}_{1}^{l} \in \mathbb{R}^{(C/r) \times C}$
and $\mathbf{W}_{2}^{l} \in \mathbb{R}^{C \times (C/r)}$ are
learned weight matrices and $\sigma$ is the sigmoid function.
The attention vector $\mathbf{a}^{l} \in [0,1]^{C}$ is
broadcast across the spatial dimensions and applied via
channel-wise multiplication:
\begin{equation}
    \hat{\mathbf{N}}_{l} = \mathbf{a}^{l} \odot \mathbf{N}_{l}
    \label{eq:lca_apply}
\end{equation}

We set $r = 16$ throughout, which means the bottleneck
introduces only $2C^{2}/16 = C^{2}/8$ additional parameters per
level---fewer than $0.02$~M parameters across all four pyramid
levels combined for $C = 256$. Despite this modest footprint,
channel attention provides a consistent mAP lift of
$0.5$--$0.8$~points in our ablation studies, with the largest
gains observed on the medium-object subset ($32^{2}$ to
$96^{2}$ pixels) of MS-COCO, where channel selectivity helps
the network resolve partially occluded instances whose
discriminative cues are concentrated in a narrow subset of
feature channels.

\subsubsection{Transformer-Style Cross-Scale Attention}

Beyond per-level channel recalibration, FusionDet optionally
supports a cross-scale attention pass that allows tokens at
one pyramid level to attend to tokens at adjacent levels.
Concretely, for each spatial position $i$ in
$\hat{\mathbf{N}}_{l}$, we form a query vector
$\mathbf{q}_{i}^{l}$ and gather key--value pairs from the
spatially corresponding $2\times2$ neighbourhood in both
$\hat{\mathbf{N}}_{l-1}$ and $\hat{\mathbf{N}}_{l+1}$:
\begin{equation}
    \text{Attn}_{i}^{l} = \text{softmax}\!\!\left(
        \frac{\mathbf{q}_{i}^{l}\,
              \mathbf{K}_{\mathcal{N}(i)}^{\top}}
             {\sqrt{d_{k}}}
    \right) \mathbf{V}_{\mathcal{N}(i)}
    \label{eq:cross_attn}
\end{equation}
where $\mathcal{N}(i)$ denotes the local neighbourhood set and
$d_{k}$ is the key dimension. Because the neighbourhood is
fixed at $2\times2$ per adjacent level, each query attends to
at most eight key--value pairs, preserving the linear
complexity guarantee established in the backbone's transformer
branch \cite{katharopoulos2020transformers}. This cross-scale
window constrains the attention span sufficiently to avoid
the quadratic blowup of full multi-scale attention, while
still enabling the detector to propagate contextual cues---such
as the co-occurrence of a traffic light and a
crosswalk---across resolution levels that would otherwise
interact only through the additive connections of the PANet
pathways.
